\section{Reprezentacja wiedzy}

\subsection{Metoda rezolucji w rachunku predykatów}

Rezolucja jest regułą wnioskowania i zarazem podstawą automatycznego dowodzenia twierdzeń. Jej główna idea polega na sprowadzeniu dowodu $T\models \beta$ do wykazania sprzeczności zbioru $T\cup\{\neg\beta\}$. Jeśli założenia razem z negacją tezy są niespełnialne, to teza wynika z założeń.

\subsubsection{Idea dowodu przez sprzeczność}

Dla teorii $T$ i zdania $\beta$ zachodzi:
\[
T\models \beta \quad\Longleftrightarrow\quad T\cup\{\neg\beta\}\text{ jest niespełnialny}.
\]
Rezolucja próbuje więc wyprowadzić klauzulę pustą, oznaczaną często $\square$. Klauzula pusta reprezentuje fałsz, czyli sprzeczność.

W rachunku zdań metoda jest decyzyjna: jeśli zbiór klauzul jest niespełnialny, rezolucja może znaleźć sprzeczność; jeśli jest spełnialny, pełne przeszukanie może to wykazać. W rachunku predykatów pierwszego rzędu sytuacja jest trudniejsza: logika pierwszego rzędu jest półrozstrzygalna. Jeśli wniosek rzeczywiście wynika, procedura może znaleźć dowód; jeśli nie wynika, przeszukiwanie może nie zakończyć się.

\subsubsection{Literały i klauzule}

Atom to formuła postaci $P(t_1,\dots,t_k)$, gdzie $P$ jest predykatem, a $t_i$ są termami. Literał to atom albo negacja atomu. Klauzula jest alternatywą literałów:
\[
L_1\vee L_2\vee\dots\vee L_m.
\]
Zbiór klauzul interpretuje się jako koniunkcję wszystkich klauzul. Klauzula pusta $\square$ jest alternatywą zera literałów, czyli jest zawsze fałszywa.

\subsubsection{Reguła rezolucji w rachunku zdań}

Jeśli mamy dwie klauzule zawierające literały komplementarne:
\[
\lambda\vee A,\qquad \neg\lambda\vee B,
\]
to możemy wyprowadzić rezolwentę:
\[
A\vee B.
\]
Reguła jest poprawna semantycznie, bo jeśli pierwsza i druga klauzula są prawdziwe, to po usunięciu pary komplementarnej pozostała alternatywa musi być prawdziwa.

Przykład:
\[
P\vee Q,\qquad \neg P\vee R
\]
dają:
\[
Q\vee R.
\]

\subsubsection{Sprowadzenie formuł do postaci klauzulowej}

Aby stosować rezolucję, formuły trzeba zamienić na postać normalną. Dla rachunku zdań używa się CNF:
\[
\bigwedge_i \bigvee_j L_{ij}.
\]
Algorytm:
\begin{enumerate}
  \item usunąć implikacje i równoważności:
  \[
  \alpha\to\beta \equiv \neg\alpha\vee\beta,
  \]
  \[
  \alpha\leftrightarrow\beta \equiv (\neg\alpha\vee\beta)\wedge(\neg\beta\vee\alpha);
  \]
  \item przesunąć negacje do atomów, używając praw de Morgana;
  \item zastosować rozdzielność $\vee$ względem $\wedge$;
  \item potraktować koniunkcję jako zbiór klauzul.
\end{enumerate}

\subsubsection{Dodatkowe problemy w rachunku predykatów}

W rachunku predykatów pojawiają się kwantyfikatory i zmienne. Dwa literały mogą nie być identyczne syntaktycznie, ale mogą stać się identyczne po podstawieniu. Przykład:
\[
Student(x),\qquad Student(Piotr).
\]
Nie są tym samym literałem, ale podstawienie $\{x/Piotr\}$ je uzgadnia. Do rezolucji pierwszego rzędu potrzebujemy więc:
\begin{itemize}
  \item postaci preneksowej i klauzulowej,
  \item standaryzacji nazw zmiennych,
  \item skolemizacji,
  \item unifikacji,
  \item reguły rezolucji z najogólniejszym unifikatorem.
\end{itemize}

\subsubsection{Postać preneksowa}

Formuła jest w postaci preneksowej, gdy wszystkie kwantyfikatory stoją z przodu:
\[
Q_1x_1Q_2x_2\dots Q_nx_n.\; M,
\]
gdzie $M$ jest formułą bez kwantyfikatorów, nazywaną macierzą.

Przekształcanie do postaci preneksowej obejmuje:
\begin{enumerate}
  \item usunięcie implikacji i równoważności,
  \item przesunięcie negacji do atomów:
  \[
  \neg\forall x\,\alpha \equiv \exists x\,\neg\alpha,\qquad
  \neg\exists x\,\alpha \equiv \forall x\,\neg\alpha,
  \]
  \item zmianę nazw zmiennych związanych, aby uniknąć kolizji,
  \item wyniesienie kwantyfikatorów na początek,
  \item przekształcenie macierzy do CNF.
\end{enumerate}

\subsubsection{Skolemizacja}

Skolemizacja usuwa kwantyfikatory egzystencjalne, zastępując zmienne egzystencjalne stałymi albo funkcjami Skolema. Nie zachowuje równoważności logicznej wprost, ale zachowuje spełnialność, co wystarcza dla rezolucji refutacyjnej.

Jeśli mamy:
\[
\exists y\, P(y),
\]
zastępujemy $y$ nową stałą $c$:
\[
P(c).
\]

Jeśli zmienna egzystencjalna zależy od wcześniejszych uniwersalnych:
\[
\forall x\exists y\, Loves(x,y),
\]
to $y$ zastępujemy funkcją Skolema $f(x)$:
\[
\forall x\, Loves(x,f(x)).
\]
Po skolemizacji wszystkie pozostałe zmienne są traktowane jako uniwersalnie kwantyfikowane, więc kwantyfikatory uniwersalne można pominąć w zapisie klauzul.

\subsubsection{Unifikacja}

Unifikacja znajduje podstawienie, które czyni dwa termy albo literały identycznymi. Najważniejszy jest najogólniejszy unifikator, MGU, czyli taki unifikator, z którego każdy inny unifikator może być otrzymany przez dalsze podstawienie.

Przykłady:
\[
P(x,a)\quad\text{i}\quad P(f(y),a)
\]
mają MGU:
\[
\{x/f(y)\}.
\]

Literały:
\[
Parent(x,John),\qquad Parent(Mary,y)
\]
mają MGU:
\[
\{x/Mary,\;y/John\}.
\]

W unifikacji trzeba stosować occurs check: zmienna nie może zostać podstawiona termem zawierającym tę samą zmienną, np. $x=f(x)$ nie powinno być dopuszczone w standardowej logice pierwszego rzędu, bo tworzyłoby nieskończony term.

\subsubsection{Reguła rezolucji pierwszego rzędu}

Jeśli mamy klauzule:
\[
L\vee A,\qquad \neg K\vee B,
\]
a $\theta$ jest MGU literałów $L$ i $K$, to rezolwentą jest:
\[
(A\vee B)\theta.
\]
Najpierw unifikujemy literały komplementarne, a potem usuwamy je i stosujemy podstawienie do reszty klauzuli.

Przykład:
\[
Human(x)\to Mortal(x),\qquad Human(Socrates).
\]
Po CNF:
\[
\neg Human(x)\vee Mortal(x),\qquad Human(Socrates).
\]
Unifikujemy $Human(x)$ z $Human(Socrates)$ przez $\{x/Socrates\}$ i otrzymujemy:
\[
Mortal(Socrates).
\]
Jeśli chcemy udowodnić $Mortal(Socrates)$ przez refutację, dodajemy $\neg Mortal(Socrates)$, a potem rezolucja daje klauzulę pustą.

\subsubsection{Procedura dowodzenia rezolucyjnego}

\begin{enumerate}
  \item Weź teorię $T$ i tezę $\beta$.
  \item Dodaj negację tezy: $T\cup\{\neg\beta\}$.
  \item Przekształć wszystkie formuły do postaci klauzulowej:
  usunięcie implikacji, przesunięcie negacji, standaryzacja zmiennych, preneks, skolemizacja, CNF, usunięcie kwantyfikatorów uniwersalnych.
  \item Wybieraj pary klauzul z unifikowalnymi literałami komplementarnymi.
  \item Dodawaj rezolwenty do zbioru klauzul.
  \item Jeśli pojawi się $\square$, dowód zakończony sukcesem.
\end{enumerate}

\subsubsection{Poprawność, pełność i problem strategii}

Rezolucja jest poprawna: każda wyprowadzona klauzula wynika logicznie z poprzednich. Jest też pełna refutacyjnie dla logiki pierwszego rzędu: jeśli zbiór klauzul jest niespełnialny, istnieje dowód rezolucyjny prowadzący do klauzuli pustej.

Pełność nie oznacza efektywności. Naiwne generowanie wszystkich rezolwent prowadzi do eksplozji kombinatorycznej. Dlatego stosuje się strategie:
\begin{itemize}
  \item rezolucję liniową,
  \item set of support,
  \item preferencję klauzul jednostkowych,
  \item subsumpcję,
  \item porządkowanie literałów,
  \item indeksowanie termów,
  \item ograniczenia głębokości albo heurystyki wyboru.
\end{itemize}

\subsubsection{Rezolucja SLD i Prolog}

W programowaniu logicznym, szczególnie w Prologu, używa się wyspecjalizowanej rezolucji SLD dla klauzul Horna. Klauzula Horna ma co najwyżej jeden literał pozytywny. Reguła:
\[
Head \leftarrow Body_1,\dots,Body_k
\]
odpowiada formule:
\[
Body_1\wedge\dots\wedge Body_k\to Head.
\]
Prolog rozwiązuje zapytanie przez próbę refutacji jego negacji, wybierając cele od lewej do prawej i klauzule w kolejności programu. To jest efektywne i praktyczne, ale strategia Prologa nie jest pełna dla wszystkich możliwych programów z powodu deterministycznego porządku przeszukiwania i możliwości zapętlenia.

\subsubsection{Co powiedzieć na obronie}

Rezolucja w rachunku predykatów polega na dowodzeniu przez sprzeczność: do teorii dodaje się negację tezy, przekształca wszystko do zbioru klauzul, usuwa kwantyfikatory egzystencjalne przez skolemizację, a literały dopasowuje przez unifikację. Reguła rezolucji usuwa parę literałów komplementarnych po zastosowaniu MGU. Wyprowadzenie klauzuli pustej oznacza, że teza wynika z teorii. Metoda jest poprawna i pełna refutacyjnie, ale wymaga strategii, bo liczba rezolwent może eksplodować.

\subsection{Główne problemy w zagadnieniach wnioskowania o działaniach}

Wnioskowanie o działaniach dotyczy systemów dynamicznych: mamy obiekty, ich własności oraz akcje zmieniające stan świata. Celem jest formalne odpowiadanie na pytania typu: co będzie prawdziwe po wykonaniu akcji, czy akcja jest wykonalna, czy jakiś cel da się osiągnąć sekwencją akcji, jakie skutki są konieczne, a jakie możliwe.

\subsubsection{Model dynamicznego świata}

Stan świata opisuje wartości fluentów, czyli własności mogących zmieniać się w czasie. Akcje przekształcają stany. W najprostszym modelu mamy:
\begin{itemize}
  \item zbiór fluentów $F$,
  \item zbiór akcji $Ac$,
  \item stany jako interpretacje fluentów,
  \item stan początkowy,
  \item funkcję przejścia $Res(A,\sigma)$ zwracającą możliwe stany po wykonaniu akcji $A$ w stanie $\sigma$.
\end{itemize}

Jeśli akcja jest deterministyczna, $Res(A,\sigma)$ ma jeden element. Jeśli niedeterministyczna, może mieć wiele wyników. Jeśli akcja jest niewykonalna, wynik może być pusty albo zgodnie z przyjętą semantyką może oznaczać brak efektu.

\subsubsection{Ontologia działań}

W modelowaniu działań trzeba zdecydować, jakie typy zjawisk dopuszczamy:
\begin{itemize}
  \item skutki deterministyczne albo niedeterministyczne,
  \item skutki pewne albo typowe, czyli domyślne,
  \item skutki bezpośrednie i pośrednie,
  \item skutki natychmiastowe i opóźnione,
  \item akcje jednoetapowe albo trwające w czasie,
  \item akcje sekwencyjne albo współbieżne,
  \item akcje elementarne albo złożone,
  \item pełną albo niepełną wiedzę o stanie i akcjach.
\end{itemize}

Każda decyzja upraszcza albo komplikuje semantykę. Prosty język akcji może wystarczyć dla sekwencyjnych działań deterministycznych, ale zawiedzie przy skutkach ubocznych, współbieżności albo wyjątkach.

\subsubsection{Frame problem}

Frame problem to problem reprezentowania i wnioskowania o tym, co nie zmienia się po wykonaniu akcji. W świecie z wieloma fluentami większość własności zwykle pozostaje bez zmian. Naiwne dopisywanie dla każdej akcji i każdego fluentu aksjomatu "ta akcja nie zmienia tej własności" prowadzi do ogromnej liczby aksjomatów.

Przykład: w Yale Shooting Problem załadowanie broni zmienia fluent $loaded$, ale nie powinno zmieniać faktu, że Fred żyje. Logika klasyczna bez dodatkowej zasady inercji nie pozwala automatycznie wywnioskować, że Fred nadal żyje po załadowaniu broni.

Rozwiązaniem jest zasada inercji: fluent zachowuje wartość, dopóki akcja nie wymusza zmiany albo nie zwalnia fluentu spod inercji. Formalizacje akcji często budują tę zasadę w semantykę, zamiast wypisywać wszystkie frame axioms.

\subsubsection{Ramification problem}

Ramification problem dotyczy skutków pośrednich. Akcja może bezpośrednio zmieniać jeden fluent, ale przez ograniczenia stanu wymuszać zmiany innych fluentów.

Przykład: jeśli obowiązuje reguła $walking\to alive$, a strzał powoduje $\neg alive$, to pośrednio powinno wynikać $\neg walking$. Nie chcemy ręcznie wypisywać każdego skutku pośredniego każdej akcji.

Trudność polega na tym, aby:
\begin{itemize}
  \item reprezentować ograniczenia statyczne,
  \item propagować skutki bezpośrednie przez te ograniczenia,
  \item nie generować nieintuicyjnych skutków,
  \item odróżnić skutki pośrednie od warunków wykonalności.
\end{itemize}

W slajdach pojawia się przykład, w którym próba zachowania spójności ograniczeń może prowadzić do dziwnych wniosków, np. "ożywienia" martwego indyka przy akcji zachęcania do chodzenia. To pokazuje, że sama minimalizacja zmian bywa niewystarczająca.

\subsubsection{Qualification problem}

Qualification problem dotyczy warunków, pod którymi akcja jest możliwa i przynosi oczekiwane skutki. W praktyce liczba potencjalnych warunków jest ogromna i nie da się ich wszystkich wypisać.

Klasyczny przykład: aby uruchomić samochód, trzeba mieć kluczyk, paliwo, sprawny akumulator, sprawne przewody i bardzo wiele innych rzeczy, w tym brak ziemniaka w rurze wydechowej. Nie da się w każdym modelu jawnie sprawdzać wszystkich nietypowych przeszkód.

Rozwiązaniem mogą być:
\begin{itemize}
  \item domyślne założenia normalności,
  \item logiki niemonotoniczne,
  \item preconditions z wyjątkami,
  \item probabilistyczne modele skuteczności działań,
  \item rozróżnienie skutków typowych i nietypowych.
\end{itemize}

\subsubsection{Problem skutków domyślnych}

W realnym świecie działania mają często skutki typowe, ale niepewne. Strzał z załadowanej broni typowo zabija, ale mogą istnieć wyjątki: broń się zacina, cel ma osłonę, strzał chybia. Logika klasyczna jest monotoniczna: dodanie nowej informacji nie usuwa wcześniejszych wniosków. Rozumowanie domyślne jest niemonotoniczne: wniosek "typowo $P$" może zostać wycofany po poznaniu wyjątku.

Modele działań z efektami domyślnymi muszą odróżniać:
\begin{itemize}
  \item skutki pewne,
  \item skutki normalne albo preferowane,
  \item skutki atypowe,
  \item minimalizację odstępstw od normalności.
\end{itemize}

\subsubsection{Niedeterminizm}

Akcja może mieć wiele możliwych rezultatów. Przykład: rzut kostką, próba otwarcia uszkodzonych drzwi, ruch robota po śliskiej powierzchni. Wtedy pytania dzielą się na:
\begin{itemize}
  \item czy cel jest konieczny po akcji, czyli zachodzi we wszystkich możliwych wynikach,
  \item czy cel jest możliwy, czyli zachodzi w co najmniej jednym wyniku,
  \item czy istnieje plan gwarantujący cel,
  \item czy istnieje strategia reagująca na obserwacje.
\end{itemize}

Niedeterminizm zwiększa złożoność planowania, bo pojedyncza sekwencja akcji może nie wystarczyć; potrzebna może być strategia warunkowa.

\subsubsection{Współbieżność}

Wiele akcji może zachodzić jednocześnie. Pojawiają się pytania:
\begin{itemize}
  \item kiedy akcje są konfliktowe,
  \item czy skutki akcji złożonej są sumą skutków akcji elementarnych,
  \item czy jedna akcja blokuje drugą,
  \item czy akcje razem mają skutek inny niż osobno.
\end{itemize}

Przykład: dwie niezależne akcje, jak otwarcie drzwi i podniesienie miski, mogą mieć łączny skutek będący koniunkcją skutków. Ale dwie akcje podnoszenia miski jedną ręką mogą osobno rozlać zupę, a razem już nie. To łamie prostą zasadę dziedziczenia skutków.

\subsubsection{Czas, akcje trwające i scenariusze}

Jeśli akcje trwają w czasie, trzeba reprezentować:
\begin{itemize}
  \item moment rozpoczęcia i zakończenia,
  \item wartości fluentów w trakcie akcji,
  \item skutki opóźnione,
  \item akcje wyzwalane przez stany,
  \item akcje wyzwalane przez inne akcje,
  \item scenariusze jako zbiory wystąpień akcji w czasie.
\end{itemize}

W modelach z czasem liniowym stan w chwili $t+1$ zależy od stanu w chwili $t$, wykonanych akcji i ewentualnych reguł dynamicznych. W modelach z czasem rozgałęzionym analizuje się wiele możliwych przyszłości.

\subsubsection{Niepełna wiedza i obserwacje}

Agent często nie zna pełnego stanu świata. Wtedy zamiast pojedynczego stanu ma zbiór stanów możliwych albo przekonanie probabilistyczne. Akcje mogą być:
\begin{itemize}
  \item ontic, czyli zmieniać świat,
  \item sensing, czyli dostarczać informacji,
  \item announcement, czyli zmieniać wiedzę agentów.
\end{itemize}

Wnioskowanie musi wtedy rozróżniać prawdę w świecie od wiedzy albo przekonań agenta. Stąd powiązanie z logikami epistemicznymi, modalnymi i dynamicznymi.

\subsubsection{Planowanie}

Wnioskowanie o działaniach jest ściśle związane z planowaniem. Typowe zapytania:
\begin{itemize}
  \item Czy po sekwencji akcji $\langle A_1,\dots,A_n\rangle$ zachodzi $\alpha$?
  \item Czy istnieje sekwencja akcji prowadząca do celu?
  \item Czy każda realizacja akcji prowadzi do celu?
  \item Czy akcja jest wykonalna w danym stanie?
  \item Jakie fluenty są obserwowalne po akcji?
\end{itemize}

W modelu niedeterministycznym albo niepełnoinformacyjnym plan może być warunkowy: następna akcja zależy od obserwacji.

\subsubsection{Co powiedzieć na obronie}

Główne problemy wnioskowania o działaniach to frame problem, ramification problem i qualification problem. Frame problem dotyczy tego, co pozostaje bez zmian; ramification problem skutków pośrednich; qualification problem ogromnej liczby warunków potrzebnych do skutecznego wykonania akcji. Dodatkowo ważne są niedeterminizm, skutki domyślne, współbieżność, czas trwania akcji, niepełna wiedza i planowanie. Formalne języki akcji próbują tak dobrać semantykę, aby uzyskać intuicyjne wnioski bez ręcznego wypisywania wszystkich wyjątków i skutków.
